\chapter{2020年江苏卷}

\section{填空题}

\begin{problem}
已知集合 \(A=\{-1,0,1,2\}\)，\(B=\{0,2,3\}\)，则 \(A\cap B=\)\fillinblank{}.
\end{problem}

\begin{problem}
已知 \(\i\) 是虚数单位，则复数 \(z=(1+\i)(2-\i)\) 的实部是\fillinblank{}.
\end{problem}

\begin{problem}
已知一组数据 \(4\)，\(2a\)，\(3-a\)，\(5\)，\(6\) 的平均数为 \(4\)，则 \(a\) 的值是\fillinblank{}.
\end{problem}

\begin{problem}
将一颗质地均匀的正方体骰子先后抛掷 \(2\) 次，观察向上的点数，则点数和为 \(5\) 的概率是\fillinblank{}.
\end{problem}

\begin{problem}
如图是一个算法流程图，若输出 \(y\) 的值为 \(-2\)，则输入 \(x\) 的值是

\bitmapfigure[max width=0.42\linewidth]{img/2020/jiangsu/q5_fig1.png}

\fillinblank{}.
\end{problem}

\begin{problem}
在平面直角坐标系 \(xOy\) 中，若双曲线 \(\dfrac{x^2}{a^2}-\dfrac{y^2}{5}=1\)（\(a>0\)）的一条渐近线方程为 \(y=\dfrac{\sqrt5}{2}x\)，则该双曲线的离心率是\fillinblank{}.
\end{problem}

\begin{problem}
已知 \(y=f(x)\) 是奇函数，当 \(x\ge0\) 时，\(f(x)=x^{\frac23}\)，则 \(f(-8)\) 的值是\fillinblank{}.
\end{problem}

\begin{problem}
已知 \(\sin^2\!\left(\dfrac\pi4+\alpha\right)=\dfrac23\)，则 \(\sin2\alpha\) 的值是\fillinblank{}.
\end{problem}

\begin{problem}
如图，六角螺帽毛坯是由一个正六棱柱挖去一个圆柱所构成的. 已知螺帽的底面正六边形边长为 \(2\mathrm{~cm}\)，高为 \(2\mathrm{~cm}\)，内孔半径为 \(0.5\mathrm{~cm}\)，则此六角螺帽毛坯的体积是

\bitmapfigure[max width=0.42\linewidth]{img/2020/jiangsu/q9_fig1.png}

\fillinblank{}\(\mathrm{cm}^3\).
\end{problem}

\begin{problem}
将函数 \(y=3\sin\!\left(2x+\dfrac\pi4\right)\) 的图象向右平移 \(\dfrac\pi6\) 个单位长度，则平移后的图象中与 \(y\) 轴最近的对称轴的方程是\fillinblank{}.
\end{problem}

\begin{problem}
设 \(\{a_n\}\) 是公差为 \(d\) 的等差数列，\(\{b_n\}\) 是公比为 \(q\) 的等比数列. 已知数列 \(\{a_n+b_n\}\) 的前 \(n\) 项和 \(S_n=n^2-n+2^n-1\)（\(n\in\N^*\)），则 \(d+q\) 的值是\fillinblank{}.
\end{problem}

\begin{problem}
已知 \(5x^2y^2+y^4=1\)（\(x\)，\(y\in\R\)），则 \(x^2+y^2\) 的最小值是\fillinblank{}.
\end{problem}

\begin{problem}
在 \(\triangle ABC\) 中，\(AB=4\)，\(AC=3\)，\(\angle BAC=90^\circ\). \(D\) 在边 \(BC\) 上，延长 \(AD\) 到 \(P\)，使得 \(AP=9\). 若 \(\overrightarrow{PA}=m\overrightarrow{PB}+\left(\dfrac32-m\right)\overrightarrow{PC}\)（\(m\) 为常数），则 \(CD\) 的长度是

\bitmapfigure[max width=0.42\linewidth]{img/2020/jiangsu/q13_fig1.png}

\fillinblank{}.
\end{problem}

\begin{problem}
在平面直角坐标系 \(xOy\) 中，已知 \(P\!\left(\dfrac{\sqrt3}{2},0\right)\)，\(A\)，\(B\) 是圆 \(C:x^2+\left(y-\dfrac12\right)^2=36\) 上的两个动点，满足 \(PA=PB\)，则 \(\triangle PAB\) 面积的最大值是\fillinblank{}.
\end{problem}

\section{解答题}

\begin{problem}
在三棱柱 \(ABC-A_1B_1C_1\) 中，\(AB\perp AC\)，\(B_1C\perp\) 平面 \(ABC\)，\(E\)，\(F\) 分别是 \(AC\)，\(B_1C\) 的中点.

\begin{enumerate}
\item 求证：\(EF\parallel\) 平面 \(AB_1C_1\)；
\item 求证：平面 \(AB_1C\perp\) 平面 \(ABB_1\).
\end{enumerate}

\bitmapfigure[max width=0.42\linewidth]{img/2020/jiangsu/q15_fig1.png}
\end{problem}

\begin{problem}
在 \(\triangle ABC\) 中，角 \(A\)，\(B\)，\(C\) 的对边分别为 \(a\)，\(b\)，\(c\)，已知 \(a=3\)，\(c=\sqrt2\)，\(B=45^\circ\).

\begin{enumerate}
\item 求 \(\sin C\) 的值；
\item 在边 \(BC\) 上取一点 \(D\)，使得 \(\cos\angle ADC=-\dfrac45\)，求 \(\tan\angle DAC\) 的值.
\end{enumerate}

\bitmapfigure[max width=0.42\linewidth]{img/2020/jiangsu/q16_fig1.png}
\end{problem}

\begin{problem}
某地准备在山谷中建一座桥梁，桥址位置的竖直截面图如图所示：谷底 \(O\) 在水平线 \(MN\) 上，桥 \(AB\) 与 \(MN\) 平行，\(OO'\) 为铅垂线（\(O'\) 在 \(AB\) 上）. 经测量，左侧曲线 \(AO\) 上任一点 \(D\) 到 \(MN\) 的距离 \(h_1\)（\(\mathrm{~m}\)）与 \(D\) 到 \(OO'\) 的距离 \(a\)（\(\mathrm{~m}\)）之间满足关系式 \(h_1=\dfrac{1}{40}a^2\)；右侧曲线 \(BO\) 上任一点 \(F\) 到 \(MN\) 的距离 \(h_2\)（\(\mathrm{~m}\)）与 \(F\) 到 \(OO'\) 的距离 \(b\)（\(\mathrm{~m}\)）之间满足关系式 \(h_2=-\dfrac{1}{800}b^3+6b\). 已知点 \(B\) 到 \(OO'\) 的距离为 \(40\mathrm{~m}\).

\begin{enumerate}
\item 求桥 \(AB\) 的长度；
\item 计划在谷底两侧建造平行于 \(OO'\) 的桥墩 \(CD\) 和 \(EF\)，且 \(CE\) 为 \(80\mathrm{~m}\)，其中 \(C\)，\(E\) 在 \(AB\) 上（不包括端点）. 桥墩 \(EF\) 每米造价 \(k\)（万元），桥墩 \(CD\) 每米造价 \(\dfrac32k\)（万元）（\(k>0\)）. 问 \(O'E\) 为多少米时，桥墩 \(CD\) 与 \(EF\) 的总造价最低？
\end{enumerate}

\bitmapfigure[max width=0.42\linewidth]{img/2020/jiangsu/q17_fig1.png}
\end{problem}

\begin{problem}
在平面直角坐标系 \(xOy\) 中，已知椭圆 \(E:\dfrac{x^2}{4}+\dfrac{y^2}{3}=1\) 的左、右焦点分别为 \(F_1\)，\(F_2\)，点 \(A\) 在椭圆 \(E\) 上且在第一象限内，\(AF_2\perp F_1F_2\)，直线 \(AF_1\) 与椭圆 \(E\) 相交于另一点 \(B\).

\begin{enumerate}
\item 求 \(\triangle AF_1F_2\) 的周长；
\item 在 \(x\) 轴上任取一点 \(P\)，直线 \(AP\) 与椭圆 \(E\) 的右准线相交于点 \(Q\)，求 \(\overrightarrow{OP}\cdot\overrightarrow{QP}\) 的最小值；
\item 设点 \(M\) 在椭圆 \(E\) 上，记 \(\triangle OAB\) 与 \(\triangle MAB\) 的面积分别为 \(S_1\)，\(S_2\)，若 \(S_2=3S_1\)，求点 \(M\) 的坐标.
\end{enumerate}

\bitmapfigure[max width=0.42\linewidth]{img/2020/jiangsu/q18_fig1.png}
\end{problem}

\begin{problem}
已知关于 \(x\) 的函数 \(y=f(x)\)，\(y=g(x)\) 与 \(h(x)=kx+b\)（\(k\)，\(b\in\R\)）在区间 \(D\) 上恒有 \(f(x)\ge h(x)\ge g(x)\).

\begin{enumerate}
\item 若 \(f(x)=x^2+2x\)，\(g(x)=-x^2+2x\)，\(D=(-\infty,+\infty)\)，求 \(h(x)\) 的表达式；
\item 若 \(f(x)=x^2-x+1\)，\(g(x)=k\ln x\)，\(h(x)=kx-k\)，\(D=(0,+\infty)\)，求 \(k\) 的取值范围；
\item 若 \(f(x)=x^4-2x^2\)，\(g(x)=4x^2-8\)，\(h(x)=4(t^3-t)x-3t^4+2t^2\)，\(0<|t|\le\sqrt2\)，\(D=[m,n]\subseteq[-\sqrt2,\sqrt2]\)，证明：\(n-m\le\sqrt7\).
\end{enumerate}
\end{problem}

\begin{problem}
已知数列 \(\{a_n\}\)（\(n\in\N^*\)）的首项 \(a_1=1\)，前 \(n\) 项和为 \(S_n\). 设 \(\lambda\) 与 \(k\) 是常数，若对一切正整数 \(n\)，均有 \(S_{n+1}^{1/k}-S_n^{1/k}=\lambda a_{n+1}^{1/k}\) 成立，则称此数列为“\(\lambda\sim k\)”数列.

\begin{enumerate}
\item 若等差数列 \(\{a_n\}\) 是“\(\lambda\sim1\)”数列，求 \(\lambda\) 的值；
\item 若数列 \(\{a_n\}\) 是“\(\dfrac{\sqrt3}{3}\sim2\)”数列，且 \(a_n>0\)，求数列 \(\{a_n\}\) 的通项公式；
\item 对于给定的 \(\lambda\)，是否存在三个不同的数列 \(\{a_n\}\) 为“\(\lambda\sim3\)”数列，且 \(a_n\ge0\)？若存在，求 \(\lambda\) 的取值范围；若不存在，说明理由.
\end{enumerate}
\end{problem}

\section{附加题：解答题}

\begin{problem}[A]
平面上点 \(A(2,-1)\) 在矩阵 \(M=\begin{bmatrix}a&1\\-1&b\end{bmatrix}\) 对应的变换作用下得到点 \(B(3,-4)\).

\begin{enumerate}
\item 求实数 \(a\)，\(b\) 的值；
\item 求矩阵 \(M\) 的逆矩阵 \(M^{-1}\).
\end{enumerate}
\end{problem}

\reuseproblemnumber
\begin{problem}[B]
在极坐标系中，已知点 \(A\!\left(\rho_1,\dfrac\pi3\right)\) 在直线 \(l:\rho\cos\theta=2\) 上，点 \(B\!\left(\rho_2,\dfrac\pi6\right)\) 在圆 \(C:\rho=4\sin\theta\) 上（其中 \(\rho\ge0\)，\(0\le\theta<2\pi\)）.

\begin{enumerate}
\item 求 \(\rho_1\)，\(\rho_2\) 的值；
\item 求出直线 \(l\) 与圆 \(C\) 的公共点的极坐标.
\end{enumerate}
\end{problem}

\reuseproblemnumber
\begin{problem}[C]
设 \(x\in\R\)，解不等式 \(2\lvert x+1\rvert+\lvert x\rvert\le4\).
\end{problem}

\begin{problem}
在三棱锥 \(A-BCD\) 中，已知 \(CB=CD=\sqrt5\)，\(BD=2\)，\(O\) 为 \(BD\) 的中点，\(AO\perp\) 平面 \(BCD\)，\(AO=2\)，\(E\) 为 \(AC\) 的中点.

\begin{enumerate}
\item 求直线 \(AB\) 与 \(DE\) 所成角的余弦值；
\item 若点 \(F\) 在 \(BC\) 上，满足 \(BF=\dfrac14BC\)，设二面角 \(F-DE-C\) 的大小为 \(\theta\)，求 \(\sin\theta\) 的值.
\end{enumerate}

\bitmapfigure[max width=0.42\linewidth]{img/2020/jiangsu/q24_fig1.png}
\end{problem}

\begin{problem}
甲口袋中装有 \(2\) 个黑球和 \(1\) 个白球，乙口袋中装有 \(3\) 个白球. 现从甲，乙两口袋中各任取一个球交换放入另一口袋，重复 \(n\) 次这样的操作，记甲口袋中黑球个数为 \(X_n\)，恰有 \(2\) 个黑球的概率为 \(p_n\)，恰有 \(1\) 个黑球的概率为 \(q_n\).

\begin{enumerate}
\item 求 \(p_1\cdot q_1\) 和 \(p_2\cdot q_2\)；
\item 求 \(2p_n+q_n\) 与 \(2p_{n-1}+q_{n-1}\) 的递推关系式和 \(X_n\) 的数学期望 \(E(X_n)\)（用 \(n\) 表示）.
\end{enumerate}
\end{problem}

